\documentclass[12pt]{article}
\usepackage[utf8]{inputenc}
\usepackage[russian]{babel}
\usepackage{graphicx}
\usepackage{array}
\usepackage[a4paper, left=1cm, right=1cm, top=1.27cm, bottom=1.27cm]{geometry}
\usepackage{setspace}
\usepackage{tabularx}
\usepackage{makecell}
\usepackage{longtable}
\usepackage{multirow}
\usepackage{caption}
\usepackage{amsmath}
\onehalfspacing


\begin{document}


\begin{table}[h]
    \centering
    \renewcommand{\arraystretch}{1.5} % Увеличиваем межстрочный интервал
    \begin{tabular}{|m{11cm}|m{6cm}|} % Первая колонка шире, вторая для лого
        \hline
        \begin{center}
            \textbf{Университет ИТМО}\\
            \textbf{Физико-технический мегафакультет}\\
            \textbf{Физический факультет}
        \end{center} 
        & 
        \begin{center}
            
\includegraphics[height=1.3cm]{logo.jpg}
        \end{center} \\
        \hline
    \end{tabular}
\end{table}

\noindent\makebox[\linewidth]{\rule{\textwidth}{2pt}}

\vspace{10pt}

\begin{table}[h]
    \centering
    \renewcommand{\arraystretch}{1.7}
    \begin{tabular}{|p{4cm}p{4cm}|p{8cm}|}
        \hline
        Группа N3150 & & К работе допущен \underline{\hspace{4cm}} \\
        \hline
        Студенты & \hfill \multirow[0pt]{3}{*}{\makecell[c]{Пахомов В.Ю \\ Фармахини Ф.Р. \\ Петров Д.О}}  & Работа выполнена \underline{\hspace{4cm}} \\
        & & \\
        & & \\
        \hline
        \multicolumn{2}{|c|}{Преподаватель \underline{\hspace{5cm}}}  & Отчет принят \underline{\hspace{5cm}} \\
        \hline
    \end{tabular}
\end{table}

\begin{center}
    \fontsize{20pt}{21.5pt} \selectfont
    \textbf{Рабочий протокол и отчет по}\\
    \textbf{лабораторной работе № 2.06}
\end{center}

\noindent\makebox[\linewidth]{\rule{\textwidth}{1pt}}
\noindent\makebox[\linewidth]{\rule{\textwidth}{1pt}}

\begin{flushleft}
    1. Цель работы.\\
    \hspace{0.63cm}Экспериментальное измерение удельной теплоёмкости материала образца.
\end{flushleft}

\begin{flushleft}
    2. Задачи, решаемые при выполнении работы.\\
    \hspace{0.63cm} Получить зависимость увеличения температуры образцов с течением времени при нагревани.
\end{flushleft}

\begin{flushleft}
    3. Объект исследования. \\
    Температруа, удельная теплоемкость.
\end{flushleft}

\begin{flushleft}
    4. Метод экспериментального исследования. \\
    \hspace{0.63cm}Проведение многократных измерений времени, за которое происходит повышение температуры
    камеры на 1 $ ^\circ C $ градус с 26 $ ^\circ C $ до 52 $ ^\circ C $. Сначала измерения проводятся с использованием пустой
    камеры, затем в камеру кладётся образец из латуни.
\end{flushleft}

\begin{flushleft}
    5. Рабочие формулы и исходные данные. \\
    $ \frac{d(T - T_{\text{окр}})}{dt} \approx \frac{2^\circ C}{t_{i+1} - t_{i-1}} $ \\
    $ \bar{x} = \frac{1}{n} \sum x_i \qquad \bar{y} = \frac{1}{n} \sum y_i $ \\
    $ b = \frac{\sum(x_i - \bar{x})(y_i - \bar{y})}{\sum(x_i - \bar{x})^2} \qquad a = \bar{y} - b \bar{x} $ \\
    $ \ln\left( \frac{d(T - T_{\text{окр}})}{dt} \right) = \ln\left( \frac{P}{C_0} \right) - \frac{K}{C_0} t $ \\
    $ P = IU $ \\
    $ \ln\left( \frac{d(T - T_{\text{окр}})}{dt} \right) = \ln\left( \frac{P}{C_0 + C} \right) - \frac{K}{C_0 + C} t $ \\
\end{flushleft}


\begin{flushleft}
    6. Измерительные приборы. \\
    \begin{table}[h]
        \centering
        \renewcommand{\arraystretch}{1.5}
        {\itshape
        \begin{tabular}{|c|c|c|c|c|}
            \hline
            № п/п & Наименование & Тип прибора & Используемый диапазон & Погрешность прибора\\
            \hline
            1 & ПКЦ-3 & Электронный & 0-99.99 $^\circ C$ & ±0.01 $^\circ C$ \\
            \hline
            2 & Весы & Электронный & 250 г & ±0.01 г \\
            \hline
        \end{tabular}
        }
    \end{table}
\end{flushleft}

\begin{flushleft}
    7. Схема установки. \\
    \includegraphics[height=11cm]{schema.png} \\
    \hspace{0.63cm}1) Моноблок ”Измерение теплоёмкости тел” с набором функциональных модулей. \\
    \hspace{0.63cm}2) Тепловая камера, в которой осуществляется нагрев образцов. \\
    \hspace{0.63cm}3) Кнопка включение питания моноблока. \\
    \hspace{0.63cm}4) Индикаторы значений напряжения и тока нагревательного элемента тепловой камеры. \\
    \hspace{0.63cm}5) Индикаторы значений текущей температуры (красный) и предельной температуры (зеленый). \\
    \hspace{0.63cm}6) Ручка регулятора напряжения, подаваемого на нагревательный элемент. \\
    \hspace{0.63cm}7) Индикатор подачи напряжения на нагревательный элемент. \\
    \hspace{0.63cm}8) Графический сенсорный дисплей для отображения графика температуры и контроля управления стендом. \\
    \hspace{0.63cm}9) Блок питания вентилятора, закрепленного под камерой (тумблер включения расположен на блоке питания). \\
    \hspace{0.63cm}10) Съёмная крышка тепловой камеры. \\
\end{flushleft}

\begin{flushleft}
    8. Результаты прямых измерений и их обработки. \\

    \newpage

    \begin{center}
        \textbf{Таблица 1. Измерения и обработка результатов эксперимента. Пустая камера} \\
    \end{center}
    \begin{table}[h!]
        \centering
        \begin{tabular}{|c|c|c|c|c|}
            \hline
            \multicolumn{3}{|c|}{Измеряемые величины} & \multicolumn{2}{c|}{Вычисляемые величины} \\
            \hline
            $T,\ ^\circ C$ & $T - T_{\text{окр}},\ ^\circ C$ & $t,\ \text{с}$ & $\frac{d(T - T_{\text{окр}})}{dt},\ ^\circ C/\text{с}$ & $\ln\left(\frac{d(T - T_{\text{окр}})}{dt}\right)$ \\
            \hline
            23.7	&	0.7	&	0	&	-	&	 -	\\
            \hline
            24.7	&	1.7	&	35	&	0.028571429	&	-3.555348061	\\
            \hline
            25.7	&	2.7	&	70	&	0.027027027	&	-3.610917913	\\
            \hline
            26.7	&	3.7	&	109	&	0.024390244	&	-3.713572067	\\
            \hline
            27.7	&	4.7	&	152	&	0.021978022	&	-3.817712326	\\
            \hline
            28.7	&	5.7	&	200	&	0.020618557	&	-3.881563798	\\
            \hline
            29.7	&	6.7	&	249	&	0.020408163	&	-3.891820298	\\
            \hline
            30.7	&	7.7	&	298	&	0.019047619	&	-3.96081317	\\
            \hline
            31.7	&	8.7	&	354	&	0.016949153	&	-4.077537444	\\
            \hline
            32.7	&	9.7	&	416	&	0.015873016	&	-4.143134726	\\
            \hline
            33.7	&	10.7	&	480	&	0.014925373	&	-4.204692619	\\
            \hline
            34.7	&	11.7	&	550	&	0.013605442	&	-4.297285406	\\
            \hline
            35.7	&	12.7	&	627	&	0.012658228	&	-4.369447852	\\
            \hline
            36.7	&	13.7	&	708	&	0.012121212	&	-4.412798293	\\
            \hline
            37.7	&	14.7	&	792	&	0.011299435	&	-4.483002552	\\
            \hline
            38.7	&	15.7	&	885	&	0.010695187	&	-4.537961436	\\
            \hline
            39.7	&	16.7	&	979	&	0.010204082	&	-4.584967479	\\
            \hline
            40.7	&	17.7	&	1081	&	0.008733624	&	-4.740574823	\\
            \hline
            41.7	&	18.7	&	1208	&	0.007575758	&	-4.882801923	\\
            \hline
            42.7	&	19.7	&	1345	&	0.006557377	&	-5.027164596	\\
            \hline
            43.7	&	20.7	&	1513	&	0.005063291	&	-5.285738584	\\
            \hline
            44.7	&	21.7	&	1740	&	0.004376368	&	-5.43153621	\\
            \hline
            45.7	&	22.7	&	1970	&	0.003539823	&	-5.643678551	\\
            \hline
            46.7	&	23.7	&	2305	&	-	&	 -	\\
            \hline
        \end{tabular}
    \end{table}

   

    \newpage

    \begin{center}
        \textbf{Таблица 2. Измерения и обработка результатов эксперимента. Камера с образцом} \\
    \end{center}
    \begin{table}[h!]
        \centering
        \begin{tabular}{|c|c|c|c|c|}
            \hline
            \multicolumn{3}{|c|}{Измеряемые величины} & \multicolumn{2}{c|}{Вычисляемые величины} \\
            \hline
            $T,\ ^\circ C$ & $T - T_{\text{окр}},\ ^\circ C$ & $t,\ \text{с}$ & $\frac{d(T - T_{\text{окр}})}{dt},\ ^\circ C/\text{с}$ & $\ln\left(\frac{d(T - T_{\text{окр}})}{dt}\right)$ \\
            \hline
            24.7	&	1.7	&	0	&	-	&	-	\\
            \hline
            25.7	&	2.7	&	54	&	0.018181818	&	-4.007333185	\\
            \hline
            26.7	&	3.7	&	110	&	0.017699115	&	-4.034240638	\\
            \hline
            27.7	&	4.7	&	167	&	0.016949153	&	-4.077537444	\\
            \hline
            28.7	&	5.7	&	228	&	0.016129032	&	-4.127134385	\\
            \hline
            29.7	&	6.7	&	291	&	0.015384615	&	-4.17438727	\\
            \hline
            30.7	&	7.7	&	358	&	0.016	&	-4.135166557	\\
            \hline
            31.7	&	8.7	&	416	&	0.015873016	&	-4.143134726	\\
            \hline
            32.7	&	9.7	&	484	&	0.013605442	&	-4.297285406	\\
            \hline
            33.7	&	10.7	&	563	&	0.01242236	&	-4.388257184	\\
            \hline
            34.7	&	11.7	&	645	&	0.012820513	&	-4.356708827	\\
            \hline
            35.7	&	12.7	&	719	&	0.012345679	&	-4.394449155	\\
            \hline
            36.7	&	13.7	&	807	&	0.010928962	&	-4.516338972	\\
            \hline
            37.7	&	14.7	&	902	&	0.00990099	&	-4.615120517	\\
            \hline
            38.7	&	15.7	&	1009	&	0.00896861	&	-4.714024591	\\
            \hline
            39.7	&	16.7	&	1125	&	0.008403361	&	-4.779123493	\\
            \hline
            40.7	&	17.7	&	1247	&	0.007692308	&	-4.86753445	\\
            \hline
            41.7	&	18.7	&	1385	&	0.006872852	&	-4.980176087	\\
            \hline
            42.7	&	19.7	&	1538	&	0.006134969	&	-5.093750201	\\
            \hline
            43.7	&	20.7	&	1711	&	0.005194805	&	-5.260096154	\\
            \hline
            44.7	&	21.7	&	1923	&	0.004347826	&	-5.438079309	\\
            \hline
            45.7	&	22.7	&	2171	&	0.004098361	&	-5.497168225	\\
            \hline
            46.7	&	23.7	&	2411	&	-	&	-	\\
            \hline
        \end{tabular}
    \end{table}
\end{flushleft}

\begin{flushleft}
    9. Расчет результатов косвенных измерений (таблицы, примеры расчетов). \\
    U = 7 в, I = 0.54 А \\
    Мощность нагревателя: $ P = UI = 3.78 $ Вт \\
    $ P - K(T - T_\text{окр}) = P \ exp(-\frac{Kt}{C + C_0}) $ \\
    Теплоемкость пустой камеры: $ C_0 = 144.6 \pm 2.4 \ \text{Дж}/^\circ C $ \\
    $ \ln\left( \frac{d(T - T_{\text{окр}})}{dt} \right) = \ln\left( \frac{P}{C_0 + C} \right) - \frac{K}{C_0 + C} t $ \\
    Теплоемкость камеры с образцом: $ C_0 + C = 158 \pm 2.5 \ \text{Дж} / ^\circ C $ \\
    $ C = 13.4 $ Дж/ $^\circ c$ \\
    Масса образца: $ m = 44.3 \ \text{г} $ \\
    $ c = \frac{C}{m} \approx 302.9 $ Дж/(кг $ \cdot ^\circ C $) \\
    По угловым коэффициентам: \\
    $ \frac{C_0 + C}{C_0} \approx 1.5063 \pm 0.064 $ \\
    $ \delta \frac{B}{B'} = \frac{B}{B'} \cdot \sqrt{(\frac{\delta B}{B})^2 + (\frac{\delta B'}{B'})^2}  \approx 0.0642 $ \\
    По свободным коэффициентам: \\
    $ \frac{C_0 + C}{C_0} \approx 1.393 \pm 0.0309 $ \\
    $ \delta (\frac{C_0 + C}{C_0}) \approx 0.0309 $ \\

\end{flushleft}

\begin{flushleft}
    10. Расчет погрешностей измерений (для прямых и косвенных измерений). \\
    $ C_0 = 144.6 \pm 2.4 \ \text{Дж}/^\circ C $ \\
    $ \sigma_C \approx 74.9 $ Дж/(кг $ \cdot ^\circ C $) \\
    $ \delta (\frac{C_0 + C}{C_0}) \approx 0.0309 $ \\
\end{flushleft}

% \newpage

\begin{flushleft}
    11. Графики (перечень графиков, которые составляют Приложение 2). \\
    \center{Пустая камера} \\
    \includegraphics[height=10cm]{chart(1).png} \\
    Камера с образцом \\
    \includegraphics[height=10cm]{chart(2).png} \\
\end{flushleft}

\begin{flushleft}
    12. Окончательные результаты. \\
    Мощность нагревателя: $ P = UI = 3.78 $ Вт \\
    Теплоемкость пустой камеры: $ C_0 = 144.6 \pm 2.4 \ \text{Дж}/^\circ C $ \\
    Теплоемкость камеры с образцом: $ C_0 + C = 158 \pm 2.5 \ \text{Дж} / ^\circ C $ \\
    Удельная теплоемкость материала: $ c = \frac{C}{m} \approx 302.9 $ Дж/(кг $ \cdot ^\circ C $) \\
\end{flushleft}

\begin{flushleft}
    13. Выводы и анализ результатов работы. \\
    Оба графика логарифмической скорости нагрева являются убывающими прямыми, что подтверждает корректность модели и линейность зависимости.
    График с образцом имеет меньший наклон, что указывает на большее значение общей теплоёмкости системы при наличии образца.
    В ходе работы была успешно определена удельная теплоёмкость образца из латуни с помощью косвенного метода, основанного на анализе графика зависимости температуры от времени.
Результаты считаются достоверными, а методика — эффективной для определения теплофизических характеристик материалов в лабораторных условиях.
Удельная теплоемкость образца оказалась максимально приближенной к табличному значению. \\

Эксперимент подтвердил применимость данного метода в учебной и исследовательской практике.

\end{flushleft}

\begin{flushleft}
    14. Дополнительные задания. \\
\end{flushleft}

\begin{flushleft}
    15. Выполнение дополнительных заданий. \\
\end{flushleft}

\begin{flushleft}
    16. Замечания преподавателя (исправления, вызванные замечаниями преподавателя, также помещают в этот пункт). \\
\end{flushleft}

\end{document}